\begin{pastexam}[label={px:ncku-107-ee-ds-07}]{成大電機所 107 學年度 問答題第 3 題}
（問答）

\textbf{3. A Mazing Problem (16 pts)}

The rat in a maze experiment is a classical one from experimental psychology, where a map is recorded
by a two-dimensional array, $maze[i][j]$, where $1\le i\le m$ and $1\le j\le p$. A value of 1 implies a
blocked path, and a 0 means one can walk right on through. The data types needed are

\quad \textbf{struct} \textit{offsets} \{ int $a$, $b$; \}

\quad \textbf{enum} directions \{N, NW, W, SW, S, SE, E, NE\};

\quad offsets \textit{move}[8];

For example, if we are at $maze[i][j]$ and we wish to find the new position $maze[g][h]$ that is
southwest of us, then we set $g=i+$move[SW].a and h $=j+$move[SW].b (see Fig. (a) and (b)).

\begin{center}
\begin{tikzpicture}[font=\scriptsize,>={Stealth[length=2mm]}]
% (a) eight directions
\begin{scope}
  \draw (-0.75,-2.1) -- (-0.75,2.1) (0.75,-2.1) -- (0.75,2.1);
  \draw (-2.1,-0.75) -- (2.1,-0.75) (-2.1,0.75) -- (2.1,0.75);
  \node[draw,circle,fill=black!15,minimum size=5mm,inner sep=0pt] (c) at (0,0) {};
  \node at (0,-0.5) {[i][j]};
  \foreach \ang/\len in {90/1.45,270/1.35,0/0.75,180/0.75,45/1.9,135/1.9,225/1.9,315/1.9} {
    \draw[->] (c) -- (\ang:\len);
  }
  \node at (0,1.7) {[i-1][j]};
  \node at (0,-1.7) {[i+1][j]};
  \node at (-1.6,1.7) {[i-1][j-1]};
  \node at (1.6,1.7) {[i-1][j+1]};
  \node at (-1.6,-1.7) {[i+1][j-1]};
  \node at (1.6,-1.7) {[i+1][j+1]};
  \node[anchor=east] at (-0.72,0) {W [i][j-1]};
  \node[anchor=west] at (0.72,0) {[i][j+1] E};
  \node at (-2.1,2.25) {NW}; \node at (0,2.25) {N}; \node at (2.1,2.25) {NE};
  \node at (-2.1,-2.25) {SW}; \node at (0,-2.25) {S}; \node at (2.1,-2.25) {SE};
  \node[font=\small] at (0,-3.1) {(a)};
\end{scope}
% (b) move table
\begin{scope}[xshift=4.9cm,font=\scriptsize\sffamily]
  \fill[black] (-1.9,0.95) rectangle (1.9,1.9);
  \draw (-1.9,-1.95) rectangle (1.9,1.9);
  \node[white] at (-1.3,1.42) {\textbf{\textit{q}}};
  \node[white,align=center] at (0,1.42) {\textbf{\textit{move[q]}}\\\textbf{\textit{.a}}};
  \node[white,align=center] at (1.25,1.42) {\textbf{\textit{move[q]}}\\\textbf{\textit{.b}}};
  \foreach \q/\a/\b [count=\r from 0] in {N/-1/0,NE/-1/1,E/0/1,SE/1/1,S/1/0,SW/1/-1,W/0/-1,NW/-1/-1} {
    \node at (-1.3,0.65-\r*0.33) {\q};
    \node at (0,0.65-\r*0.33) {\a};
    \node at (1.25,0.65-\r*0.33) {\b};
  }
  \node[font=\small\rmfamily] at (0,-3.1) {(b)};
\end{scope}
% (c) maze
\begin{scope}[xshift=8.3cm,yshift=1.75cm,x=0.62cm,y=0.62cm,font=\small]
  \draw[step=1] (0,0) grid (5,-5);
  \foreach \row [count=\r from 0] in {{0,0,0,1,1},{0,1,1,0,1},{0,0,0,0,0},{1,1,1,0,1},{0,1,1,0,0}} {
    \node at (-0.6,-\r-0.5) {\r};
    \foreach \v [count=\c from 0] in \row { \node at (\c+0.5,-\r-0.5) {\v}; }
  }
  \foreach \c in {0,...,4} { \node at (\c+0.5,0.55) {\c}; }
  \draw[->] (6.1,0.55) -- (5.3,0.55);
  \node at (6.5,0.55) {$j$};
  \draw[->] (-0.6,-6.0) -- (-0.6,-5.2);
  \node at (-0.6,-6.5) {$i$};
  \draw[->] (-0.5,0.7) -- (0,0);
  \node[font=\scriptsize] at (-0.7,1.0) {exit};
  \draw[->] (5.7,-5.6) -- (5,-5);
  \node[font=\scriptsize] at (5.6,-6.0) {entrance};
  \node at (2.5,-7.3) {(c)};
\end{scope}
\end{tikzpicture}
\end{center}

The pseudo code is shown in the following:
\begin{minted}[fontsize=\small, linenos, numbersep=6pt, xleftmargin=2em]{text}
initialize list to the maze entrance coordinates and direction north;
while (list is not empty)
{
  (i, j, dir) = coordinates and direction from end of list;
  delete last element of list;
  while (there are more moves from (i, j)) // i.e., dir < 8
  {
     (g, h) = coordinates of next move;
     if ((g == m) && (h == p)) success;
     if ((!maze[g][h] && (!mark[g][h]))
     {
        mark[g][h] = 1;
        dir = next direction to try;
        add (i, j, dir) to end of list;
        (i, j, dir) = (g, h, N);  //move to new location
     }
  }
}
cout << "No path in maze." << endl;
\end{minted}
\begin{enumerate}[label=\alph*.]
\item \textbf{(2 pts)} What kind of the data structure that you can use it to implement the
\textit{list} in the above pseudo code? Why?
\item \textbf{(2 pts)} How do you handle the problem when a rat touches the boundary of a maze?
\item \textbf{(2 pts)} What is the purpose of the matrix \textit{mark} which appears in Line 10 and Line
12?
\item \textbf{(5 pts)} Please show the path that is found by the program in the map of Fig. (c) in term
of the index of the matrix.
\item \textbf{(2 pts)} Can you find the path if there exists at least one route from entrance to exit
according to the algorithm?
\item \textbf{(3 pts)} Assume the cost will be increased by one if you go through a grid in the maze,
and the cost of a path is the number of grids in the path. Dose the proposed algorithm find the path
with the smallest cost? If the answer is No, what kind of algorithm and the data structure that you
should use in order to find the path with the smallest cost?
\end{enumerate}
\end{pastexam}

\begin{pastsoln}
這是 Horowitz 課本的老鼠走迷宮，原題有幾處要先說清楚：
\begin{itemize}
\item 文字說下標從 $1$ 開始、出口是 $(m, p)$，但圖 (c) 的下標是 $0$ 到 $4$，入口在右下角 $(4, 4)$，出口在左上角
$(0, 0)$。 下面照圖作答：從 $(4, 4)$ 出發，第 9 行的條件讀成「$(g, h)$ 是出口 $(0, 0)$」。
\item 第 10 行原卷少一個右括號，意思是 \texttt{if (!maze[g][h] \&\& !mark[g][h])}。
\item 內層迴圈只在第 13 行改 \texttt{dir}；那一格走不通時也要換下一個方向，否則會停在原地，所以讀成「不能走
就 \texttt{dir++}」（Horowitz 的原程式有這個 else）。 入口也先設 \texttt{mark} 為 $1$，不讓老鼠走回入口。
\item 「下一個方向」照 \textbf{enum} 的順序 N、NW、W、SW、S、SE、E、NE（enum 決定 \texttt{dir} 的數值）。
圖 (b) 表格列的順序是 N、NE、E、SE、…，那只是查表，不影響 \texttt{dir} 加一後是哪個方向。 最後也附上照表格
順序的結果。
\end{itemize}

\textbf{a.} 用\textbf{堆疊}。 程式只在 \textit{list} 的尾端加入（第 14 行）和取出（第 4--5 行），後進先出。 這正是
回溯的需要：走進死路時，要退回的是\textbf{最近}經過、還有方向沒試完的那一格，也就是最後放進去的那一筆。
堆疊裡由底到頂存的正好是目前從入口走到這裡的路徑。

\textbf{b.} 在迷宮外圍加一圈 $1$（牆）：開成 $(m+2) \times (p+2)$ 的陣列，第 $0$ 列、第 $m+1$ 列、第 $0$ 行、第 $p+1$
行全設為 $1$。 這樣任何一格往八個方向走，算出來的 $(g, h)$ 都在陣列範圍內，碰到外圍的牆自然走不通，第 10 行
不必另外檢查邊界。 （另一種做法是每次檢查 $0 \le g < m$、$0 \le h < p$，超出範圍就當牆。）

\textbf{c.} \textit{mark} 記錄哪些格子已經走過（放進過路徑）。 第 10 行只走沒走過的格子，第 12 行在走進去時標記。
沒有它，老鼠會在兩格之間來回走（例如 $A \to B \to A \to B \cdots$）或在環裡繞圈，程式不會結束；有了它，每一格最多
進入一次，總時間是 $O(mp)$（每格最多試八個方向）。

\textbf{d.} 找到的路徑（列, 行）：
\[
(4,4) \to (3,3) \to (2,3) \to (1,3) \to (0,2) \to (0,1) \to (0,0).
\]
追蹤（「試」是依 enum 順序試過的方向，$\times$ 表示是牆或超出範圍）：
\begin{center}
\small
\begin{tabular}{lll}
\toprule
所在 & 依序試的方向 & 結果 \\
\midrule
$(4,4)$ & N $(3,4)$ $\times$，NW $(3,3)$ 通 & 把 $(4,4,\text{W})$ 放進堆疊，走到 $(3,3)$ \\
$(3,3)$ & N $(2,3)$ 通 & 放進 $(3,3,\text{NW})$，走到 $(2,3)$ \\
$(2,3)$ & N $(1,3)$ 通 & 放進 $(2,3,\text{NW})$，走到 $(1,3)$ \\
$(1,3)$ & N $(0,3)$ $\times$，NW $(0,2)$ 通 & 放進 $(1,3,\text{W})$，走到 $(0,2)$ \\
$(0,2)$ & N 超出，NW 超出，W $(0,1)$ 通 & 放進 $(0,2,\text{SW})$，走到 $(0,1)$ \\
$(0,1)$ & N 超出，NW 超出，W $(0,0)$ 是出口 & 成功 \\
\bottomrule
\end{tabular}
\end{center}
這次沒有走進死路，堆疊裡的五格加上目前所在 $(0,1)$ 和出口 $(0,0)$，就是上面的路徑。

若「下一個方向」照圖 (b) 表格的順序（N、NE、E、SE、S、SW、W、NW），走法會不同：$(4,4)$ 先走 W 到 $(4,3)$，
往北到 $(3,3)$、$(2,3)$、$(1,3)$；在 $(1,3)$ 先試到 SE $(2,4)$，那是死路（周圍不是牆就是走過的格子），退回 $(1,3)$
改走 SW $(2,2)$，再往西 $(2,1)$、$(2,0)$，往北 $(1,0)$，到出口 $(0,0)$。 路徑是
$(4,4)$、$(4,3)$、$(3,3)$、$(2,3)$、$(1,3)$、$(2,2)$、$(2,1)$、$(2,0)$、$(1,0)$、$(0,0)$。 作答時寫清楚採用哪一種順序。

\textbf{e.} 能。 只要入口到出口有路，演算法一定找得到：它是深度優先搜尋，每一格的八個方向都會試完才放棄，
退回去時會繼續試上一格剩下的方向。 一格只有在「從它出發的所有方向都試過了」才會永久離開堆疊，所以從
入口能走到的格子全部都會被標記、被走過（和 DFS 會拜訪起點能到的所有頂點同理，資結書第~21.4~節）。 出口
和某個能走到的格子相鄰時，從那一格試到出口的方向就成功。 只有入口根本走不到出口時，堆疊才會清空，印出
``No path in maze.''。

\textbf{f.} 不一定。 DFS 找到的是「第一條」路，不保證最短。 上面 d 找到的路經過 $7$ 格，但
$(4,4) \to (3,3) \to (2,2) \to (2,1) \to (1,0) \to (0,0)$ 只經過 $6$ 格。 要找格數最少的路，改用\textbf{廣度優先搜尋}
（BFS），資料結構用\textbf{佇列}（資結書第~21.3~節）：從入口出發，一層一層往外擴，先找到的一定是步數最少的；
每格記下從哪一格來，到出口後倒推就是最短路徑。 每一步的成本都是 $1$，所以不需要 Dijkstra；若每格成本不同
才要用 Dijkstra 演算法配優先佇列。
\end{pastsoln}
